Bodový graf závislosti poměru průměrné ke minimální mzdě (nebo mediánové mzdy k~dolnímu decilu) na \ac{KS} pokrytí ukazuje, jak silné \ac{KV} instituce stlačují vzdálenost mezi minimálním a~mediánovým standardem: v~zemích s~erga omnes je tento poměr nižší --- tedy minimální mzda (nebo kolektivní smluvní tarif) je blíže mediánu.
V~ČR je zákonná minimální mzda relativně vzdálena od mediánu --- přibližně \SI{53}{\percent} mediánu v~roce 2023, zatímco v~DK nebo BE nepotřebují zákonnou minimální mzdu vůbec, protože smluvní minimum je nastaveno kolektivně vysoko.
Tento graf demonstruje, že minimální mzda a~\ac{KS} pokrytí jsou substituty: v~zemích se silným \ac{KV} pokrytím je zákonná minimální mzda méně relevantní (\ac{KS} ji nastavuje výše), zatímco v~zemích s~nízkým pokrytím (ČR) zákonná minimální mzda je posledním záchranným sítem --- ale sítem s~hrubými oky.
Politická implikace: namísto soupeření o~výši zákonné minimální mzdy by měla být prioritou reforma \ac{KV} pokrytí, která by zákonné minimum přirozeně odsunula do role pouhé bezpečnostní sítě pro ty, kteří jsou mimo dosah \ac{KS} --- čímž by se oba nástroje staly komplementárními, nikoliv soupeřícími.

Bodový graf závislosti poměru průměrné ke minimální mzdě (nebo mediánové mzdy k~dolnímu decilu) na \ac{KS} pokrytí ukazuje, jak silné \ac{KV} instituce stlačují vzdálenost mezi minimálním a~mediánovým standardem: v~zemích s~erga omnes je tento poměr nižší --- tedy minimální mzda (nebo kolektivní smluvní tarif) je blíže mediánu.
V~ČR je zákonná minimální mzda relativně vzdálena od mediánu --- přibližně \SI{53}{\percent} mediánu v~roce 2023, zatímco v~DK nebo BE nepotřebují zákonnou minimální mzdu vůbec, protože smluvní minimum je nastaveno kolektivně vysoko.
Tento graf demonstruje, že minimální mzda a~\ac{KS} pokrytí jsou substituty: v~zemích se silným \ac{KV} pokrytím je zákonná minimální mzda méně relevantní (\ac{KS} ji nastavuje výše), zatímco v~zemích s~nízkým pokrytím (ČR) zákonná minimální mzda je posledním záchranným sítem --- ale sítem s~hrubými oky.
Politická implikace: namísto soupeření o~výši zákonné minimální mzdy by měla být prioritou reforma \ac{KV} pokrytí, která by zákonné minimum přirozeně odsunula do role pouhé bezpečnostní sítě pro ty, kteří jsou mimo dosah \ac{KS} --- čímž by se oba nástroje staly komplementárními, nikoliv soupeřícími.

Bodový graf závislosti poměru průměrné ke minimální mzdě (nebo mediánové mzdy k~dolnímu decilu) na \ac{KS} pokrytí ukazuje, jak silné \ac{KV} instituce stlačují vzdálenost mezi minimálním a~mediánovým standardem: v~zemích s~erga omnes je tento poměr nižší --- tedy minimální mzda (nebo kolektivní smluvní tarif) je blíže mediánu.
V~ČR je zákonná minimální mzda relativně vzdálena od mediánu --- přibližně \SI{53}{\percent} mediánu v~roce 2023, zatímco v~DK nebo BE nepotřebují zákonnou minimální mzdu vůbec, protože smluvní minimum je nastaveno kolektivně vysoko.
Tento graf demonstruje, že minimální mzda a~\ac{KS} pokrytí jsou substituty: v~zemích se silným \ac{KV} pokrytím je zákonná minimální mzda méně relevantní (\ac{KS} ji nastavuje výše), zatímco v~zemích s~nízkým pokrytím (ČR) zákonná minimální mzda je posledním záchranným sítem --- ale sítem s~hrubými oky.
Politická implikace: namísto soupeření o~výši zákonné minimální mzdy by měla být prioritou reforma \ac{KV} pokrytí, která by zákonné minimum přirozeně odsunula do role pouhé bezpečnostní sítě pro ty, kteří jsou mimo dosah \ac{KS} --- čímž by se oba nástroje staly komplementárními, nikoliv soupeřícími.

Bodový graf závislosti poměru průměrné ke minimální mzdě (nebo mediánové mzdy k~dolnímu decilu) na \ac{KS} pokrytí ukazuje, jak silné \ac{KV} instituce stlačují vzdálenost mezi minimálním a~mediánovým standardem: v~zemích s~erga omnes je tento poměr nižší --- tedy minimální mzda (nebo kolektivní smluvní tarif) je blíže mediánu.
V~ČR je zákonná minimální mzda relativně vzdálena od mediánu --- přibližně \SI{53}{\percent} mediánu v~roce 2023, zatímco v~DK nebo BE nepotřebují zákonnou minimální mzdu vůbec, protože smluvní minimum je nastaveno kolektivně vysoko.
Tento graf demonstruje, že minimální mzda a~\ac{KS} pokrytí jsou substituty: v~zemích se silným \ac{KV} pokrytím je zákonná minimální mzda méně relevantní (\ac{KS} ji nastavuje výše), zatímco v~zemích s~nízkým pokrytím (ČR) zákonná minimální mzda je posledním záchranným sítem --- ale sítem s~hrubými oky.
Politická implikace: namísto soupeření o~výši zákonné minimální mzdy by měla být prioritou reforma \ac{KV} pokrytí, která by zákonné minimum přirozeně odsunula do role pouhé bezpečnostní sítě pro ty, kteří jsou mimo dosah \ac{KS} --- čímž by se oba nástroje staly komplementárními, nikoliv soupeřícími.

\input{texparts/python/coverage_ratio_scatter}



